\clearpage\thispagestyle{empty}
\pdfbookmark[0]{Quellen zur Besprechung von Hardings Himmelsatlas}{harding1809-sources}
\begin{center}\textbf{Quellen zur Besprechung von Hardings Himmelsatlas}\end{center}
\noindent Diese KI-Übertragung ins heutige Deutsch folgt Gauß’ Besprechung auf S.~510--512 der \href{https://gdz.sub.uni-goettingen.de/id/PPN236007467}{\textit{Werke}, Bd.~VI} (Göttingen, 1874). Sie erschien ohne Namensnennung unter der Überschrift \emph{Göttingen} in den \textit{Göttingischen gelehrten Anzeigen}, Stück~61, 17. April 1809, S.~\href{https://images.p.sub.uni-goettingen.de/iiif/image/gdz:PPN319721507_1809:00000603/full/1750,/0/default.jpg}{601}--\href{https://images.p.sub.uni-goettingen.de/iiif/image/gdz:PPN319721507_1809:00000609/full/1750,/0/default.jpg}{607}. Das digitalisierte Exemplar stammt aus der Niedersächsischen Staats- und Universitätsbibliothek Göttingen.
\par\smallskip\noindent Der Erstdruck nennt das letzte Blatt der Aufzählung IX. Die Werkausgabe nennt es dort IV und führt zu Beginn die Blätter I, II, V und IX auf. Beide Drucke geben den Endwert 81 der Rektaszension ohne Gradzeichen an. Die Übertragung folgt der Werkausgabe.
\par\smallskip\noindent Gauß verweist auf seinen Aufsatz \hyperlink{grenzen1804.1}{\emph{Über die Grenzen der geozentrischen Positionen der Planeten}} vom August 1804. Delambres Zitat aus Plutarch bedeutet \emph{ein Werk des vom Glück begünstigten Könnens, nicht des Zufalls}. Die Wortfolge weicht von \href{https://penelope.uchicago.edu/Thayer/E/Roman/Texts/Plutarch/Lives/Timoleon*.html#36.4}{Plutarchs \textit{Timoleon}, 36.4} ab. \href{https://www.davidrumsey.com/luna/servlet/detail/RUMSEY~8~1~372247~90138986:IX-Atlas-novus-coelestis-XXVII-tabu}{Hardings Blatt IX} in der David Rumsey Historical Map Collection gehört zum Atlas von 1822 mit 27 Blättern. KI-Übertragung und KI-Satz.
\clearpage
